\documentclass[10pt,a4paper]{amsart}
\usepackage[margin=19mm]{geometry}
\usepackage{amsmath,amssymb,mathtools}
\usepackage[scr=rsfs,cal=euler]{mathalfa}
\usepackage{xfrac}
\usepackage[unicode,hidelinks,bookmarks=false]{hyperref}
\newcommand{\ie}{i.e.\ }
\newcommand{\cf}{cf.\ }
\newcommand{\resp}{respectively\ }
\newcommand{\Subsection}{\subsection}
\providecommand{\nobreakdash}{\nobreak\hbox{-}}
\setlength{\emergencystretch}{2em}
\multlinegap0pt
\usepackage{amssymb}
\usepackage{enumerate}
\usepackage{tikz}
\newtheorem{thm}{Theorem}
\newtheorem{lem}{Lemma}
\DeclareMathOperator{\Ass}{Ass}
\DeclareMathOperator{\diam}{diam}
\DeclareMathOperator{\dist}{dist}
\newcommand{\mm}{\mathfrak m}
\DeclareMathOperator{\supp}{supp}
\title{Associated primes of powers of closed neighborhood ideals and diameters of graphs}
\author{Ha Thi Thu Hien, Thanh Vu}
\date{Source: arXiv:2604.23428v2 (CC BY 4.0)}
\begin{document}
\maketitle

\noindent\emph{Source and licence.} Associated primes of powers of closed neighborhood ideals and diameters of graphs. Version arXiv:2604.23428v2 (CC BY 4.0), distributed by arXiv under the Creative Commons Attribution 4.0 International licence, http://creativecommons.org/licenses/by/4.0/, as recorded in the arXiv licence metadata for this exact version and shown on the abstract page https://arxiv.org/abs/2604.23428v2. Full licence notice: \texttt{licenses/arxiv-2604.23428-CC-BY-4.0.txt}.\par\smallskip

\noindent\textit{Editorial scope.} Counted unit: the article's thm thm\_diam (source0.tex lines 291--297). The article's complete original proof of it is delivered with it (source0.tex lines 482--534, one \texttt{proof} environment of 53 lines, which the article places away from the statement). No mathematics is written, repaired or re-derived: the statement and the proof are the article's own lines, in the article's own notation, and no retained line is substituted or re-flowed.  Retained uncounted as the source-local context the proof needs: lines 306--308; lines 410--412; lines 459--461.  Nothing was omitted from any retained span, so the package carries no omission record.  No retained line contains a citation, so the wrapper carries no bibliography and no bibliography entry.  No retained line contains a figure environment, an image-inclusion command or drawing code, so no image is delivered and the package declares no asset dependency.  Editorial formatting only, in addition: an AMS \texttt{amsart} preamble that restates the article's own declarations and macros used by the retained lines, namely the environments \texttt{thm}, \texttt{lem} on separate counters, together with the standard packages those lines need.  The retained environments print in this document as Theorem 1, Lemma 1 in order of appearance, whereas the article prints the same items with the numbering of its own full document.  The complete native source of the article as distributed in the arXiv e-print for this exact version is bundled unchanged as \texttt{sources/199165/source0.tex}.
\begin{lem}\label{lem_quotient}
Let $J$ be a squarefree monomial ideal. Then $\mm$ is an associated prime of $S/J^t$ if and only if there exists a monomial $f$ such that $\deg_i(f) < t$ for all $i$ and $J : f = \mm$.
\end{lem}

\begin{lem}\label{lem_extreme_graph}
We have $\diam(G_t) = 7t - 1$ and $\mm \in \Ass(S/NI(G_t)^{t+1})$.
\end{lem}

\begin{lem}\label{lem_disjoint_ind}
Let $G$ be a connected graph with $\diam(G) = d$. Then there exist $\left\lfloor \frac{d}{3} \right\rfloor + 1$ vertices of $G$ whose closed neighborhoods are pairwise disjoint.
\end{lem}

\begin{thm}\label{thm_diam}
Let $G$ be a simple connected graph and $t \ge 2$ be an integer. Assume that $\mathfrak{m}$ is an associated prime of $S/NI(G)^t$. Then 
\[
\operatorname{diam}(G) \le 7t - 8,
\]
and this bound is sharp.
\end{thm}

\begin{proof}[Proof of Theorem \ref{thm_diam}]
By Lemma \ref{lem_extreme_graph}, it suffices to show that if $\mm \in \Ass(S/NI(G)^t)$, then $\diam(G) \leq 7t - 8$. We therefore assume that $\mm \in \Ass(S/NI(G)^t)$. By Lemma \ref{lem_quotient}, there exists a monomial $f$ such that $\mm = NI(G)^t : f$. In particular, $f \notin NI(G)^t$. Let $U = \supp(f)$ and 
\[
W = \{u \in U \mid N_G[u] \subseteq U\}.
\]
For ease of reading, we divide the proof into several steps.

\medskip

\noindent \textbf{Step 1.} $W$ is nonempty and $N_G[W] = U$.

Clearly, $f \neq 1$, hence $U \neq \emptyset$. Let $u \in U$. Since $x_u f \in NI(G)^t$, there exist vertices $i_1, \ldots, i_t$ such that $x_{N[i_1]} \cdots x_{N[i_t]}$ divides $x_u f$. In particular, $N[i_j] \subseteq \supp(uf) = \supp(f) = U$ for all $j = 1, \ldots, t$. Hence, $i_1, \ldots, i_t \in W$, and therefore $W$ is nonempty.

Furthermore, since $f \notin NI(G)^t$, there exists some $j$ such that $u \in N[i_j]$. Thus, $u \in N_G[i_j] \subseteq N_G[W]$, and hence $N_G[W] = U$.

\medskip

\noindent \textbf{Step 2.} $N_G[U] = V(G)$.

If $U = V(G)$, there is nothing to prove. Let $v \notin U$. Then $x_v f \in NI(G)^t$, so there exist vertices $v_1, \ldots, v_t$ such that $x_{N[v_1]} \cdots x_{N[v_t]}$ divides $x_v f$. Note that this product does not divide $f$, so there exists some $j$ such that $v \in N[v_j]$.

Since $v \notin U$, we have $\deg_v(x_v f) = 1$. Hence,
\[
N_G[v_j] \setminus \{v\} \subseteq \supp(f) = U.
\]
If $v_j \neq v$, then $v_j \in U$. Hence, $v \in N_G[v_j] \subseteq N_G[U]$. If $v_j = v$, then there exists $u_j \in N_G[v] \setminus \{v\}$. Then $u_j \in U$, which also implies that $v \in N_G[u_j] \subseteq N_G[U]$. Thus, $N_G[U] = V(G)$.

\medskip

\noindent \textbf{Step 3.} Assume that $G[W]$ has $s$ connected components $W_1, \ldots, W_s$. Let $U_i = N_G[W_i]$ and $V_i = N_G[U_i]$. Then $\dist_G(V_i, V_j) \leq s - 1$ for all $i, j = 1, \ldots, s$, where by convention $\dist_G(U, \emptyset) = 0$ for any subset $U \subseteq V(G)$.

We may assume that $s \geq 2$. Since $G$ is connected, the sets $V_i$ are pairwise connected through the graph. As there are $s$ components, it follows that $\dist_G(V_i, V_j) \leq s - 1$ for all $i, j$.

\medskip

\noindent \textbf{Step 4.} Final step.

Let $\diam(W_i) = a_i$ for $i = 1, \ldots, s$. By Step 1, $\diam(U_i) \leq \diam(W_i) + 2$. By Step 2, $\diam(V_i) \leq \diam(U_i) + 2 \leq \diam(W_i) + 4$. By Step 3,
\[
\diam(G) \leq \sum_{i=1}^s \diam(W_i) + 4s + (s - 1).
\]

Assume that $\diam(W_i) = a_i$. By Lemma \ref{lem_disjoint_ind}, we can choose $\left\lfloor \frac{a_i}{3} \right\rfloor + 1$ vertices in $W_i$ whose closed neighborhoods are pairwise disjoint. Since $f \notin NI(G)^t$, there can be at most $t - 1$ such vertices in total. Therefore,
\[
\sum_{i=1}^s \left(\left\lfloor \frac{a_i}{3} \right\rfloor + 1\right) \leq t - 1.
\]

Under this constraint, the sum $\sum_{i=1}^s a_i$ is maximized when each $a_i$ is of the form $a_i = 3b_i + 2$. The condition then becomes $\sum_{i=1}^s b_i \leq t - s - 1$. Hence,
\[
\diam(G) \leq \sum_{i=1}^s a_i + 5s - 1 \leq 3(t - s - 1) + 2s + 5s - 1 = 3t + 4s - 4.
\]
Since $s \leq t - 1$, we conclude that $\diam(G) \leq 7t - 8$.
\end{proof}

\end{document}
